% H9-T24: full geometric generic triple-cover space in every characteristic.
% Run JAC1964-20261121; /root/h9_height; GPT-6 Astra Pro.

\subsection*{H9-T24. The full generic triple cover in every genus at least two}

\begin{theorem}[H9-T24]
Let $A/k$, $n=2m\ge4$, and its genuine Lefschetz pencil
$f:\widetilde A\to\mathbf P^1$ be as in H9-T19, with
critical set $\Sigma$. Fix $g\ge2$. In every characteristic,
the stack of smooth geometrically connected genus-$g$
degree-three covers of the fixed target $\mathbf P^1$,
unramified above $\Sigma$, is irreducible. A nonempty open
set in an explicit affine space of twisted binary cubic
forms dominates its dense balanced locus.

Take the geometric generic cover in this parameter space,
write $F$ for the algebraic closure of its parameter
function field, and put
\[
       K=F(t),\qquad L=F(C),\qquad [L:K]=3.
\]
All homologically trivial rational codimension-$m$ cycles
on $Y/K$ and $Y_L$ are admissible, and
\[
 h_L(z,w)=\frac13 h_K(\operatorname{cor}z,
                              \operatorname{cor}w),\qquad
 \operatorname{rad}h_L
      =\operatorname{cor}^{-1}(\operatorname{rad}h_K).
                                                               \tag{H9.135}
\]
Thus the generic triple cover adds no directions to the
height quotient. For $m=2$, the full homologically trivial
height group is positive semidefinite under 032 and H9-T8.
For higher $m$, the formula transfers the original axis
form with the indicated positive scalar.

The assertion is geometric generic, with the stated
extension of constants. It neither identifies the Chow
norm kernel with zero nor asserts the formula for every
special triple cover. The genus-one case is not included.
\end{theorem}

\begin{proof}
We first give the parameter argument, including its
characteristic-three scope. Throughout, $\operatorname{Sym}^3$
denotes the ordinary symmetric quotient of the third
tensor power, so its local basis is $x^3,x^2y,xy^2,y^3$,
without binomial coefficients.

\paragraph{Cubic algebras over an arbitrary base.}
We use the equivalence in Wood, \emph{Parametrizing quartic
algebras over an arbitrary base}, \S2, Theorem~2.1,
printed p.~1073: a finite locally free rank-three algebra
$\mathcal C$ over any scheme corresponds to a rank-two
bundle $E$ and a section
\[
       \phi\in\operatorname{Sym}^3E\otimes\det(E)^{-1},
                  \qquad \mathcal C/\mathcal O\simeq E^\vee.
                                                               \tag{H9.136}
\]
This equivalence commutes with base change, with no
restriction on the characteristic. We use this algebra
construction, so do not require a choice of a relative
Gorenstein embedding or a trace splitting by $1/3$.

For precision, Wood's normalized multiplication table
on printed p.~1074, equation~(1), associates
$(A_0x^3+A_1x^2y+A_2xy^2+A_3y^3)(x\wedge y)^\vee$
with a basis $1,\omega,\theta$ satisfying
\[
 \omega\theta=-A_0A_3,\qquad
 \omega^2=-A_0A_2+A_1\omega-A_0\theta,\qquad
 \theta^2=-A_1A_3+A_3\omega-A_2\theta.
                                                               \tag{H9.137}
\]
These identities also hold in characteristics two and
three. They will give the only splitting-type bound
needed below.

\paragraph{Every smooth connected cover has unobstructed cubic coefficients.}
Let $\pi:C\to\mathbf P^1$ be any cover in the theorem,
over any algebraically closed extension of $k$, and set
$E=(\pi_*\mathcal O_C/\mathcal O)^\vee$.
The unit is a line subbundle: its image in every fibre
of the rank-three algebra is nonzero. Split
\[
                  E=\mathcal O(e)\oplus\mathcal O(f_0),
                            \qquad e\le f_0.
\]
We temporarily use $f_0$ for the larger splitting degree,
distinct from the pencil morphism $f$.
Since $H^0(C,\mathcal O_C)=k$ and
$H^1(\mathbf P^1,\mathcal O)=0$, the unit exact sequence
gives $H^0(E^\vee)=0$. Euler characteristic in the same
sequence gives
\[
                     1\le e\le f_0,\qquad e+f_0=g+2.
                                                               \tag{H9.138}
\]
The coefficient bundle in (H9.136) is
\[
 \mathcal B_E:=\operatorname{Sym}^3E\otimes\det(E)^{-1}
  =\mathcal O(2e-f_0)\oplus\mathcal O(e)
          \oplus\mathcal O(f_0)\oplus\mathcal O(2f_0-e).
                                                               \tag{H9.139}
\]
If $f_0>2e$, its first coefficient $A_0$ must vanish.
At the generic point of $\mathbf P^1$, (H9.137) would
give $\omega(\omega-A_1)=0$ in the degree-three field
$k(C)$. Both factors are nonzero, since $1,\omega,\theta$
are linearly independent over $k(t)$. This is impossible.
Consequently
\[
                         f_0\le2e,
               \qquad H^1(\mathbf P^1,\mathcal B_E)=0.
                                                               \tag{H9.140}
\]
The second assertion follows directly from (H9.139):
all four line-bundle degrees are nonnegative. Smoothness
and connectedness of $C$ imply its integrality, as used
in the field argument; unramifiedness over the nonempty
$\Sigma$ also ensures the cover is separable.

\paragraph{The balanced parameter space.}
Set
\[
 e_g=\left\lfloor\frac{g+2}{2}\right\rfloor,
 \quad f_g=\left\lceil\frac{g+2}{2}\right\rceil,
 \quad E_g=\mathcal O(e_g)\oplus\mathcal O(f_g),
 \quad P_g=H^0(\mathbf P^1,\mathcal B_{E_g}).
                                                               \tag{H9.141}
\]
Thus $P_g$ is an affine space of dimension $2g+8$.
Apply (H9.136) to its universal section on
$\mathbf P^1\times P_g$. This gives a genuine finite
locally free degree-three cover of that product.
Let $U_g\subset P_g$ be the open set where its curves
are smooth and geometrically connected and its map is
unramified above $\Sigma$. The genus is $g$ by (H9.138).
Smoothness and connectedness are open for this proper
flat family; the unramified condition at the finite set
$\Sigma$ is open as well. We now exhibit an actual
T23 family in this space, which proves $U_g\ne\varnothing$.

Write $g=2a+3b$ with $b=0$ for even $g$ and $b=1$ for
odd $g$. Then
\[
                       e_g=a+b+1,\qquad f_g=a+2b+1.
                                                               \tag{H9.142}
\]
In characteristic different from three use the tame
family (H9.127). Homogenize its right side as a section
of $\mathcal L^3$, where $\mathcal L=\mathcal O(a+b+1)$.
The branch exponents are one or two, with exponent two
at infinity. There are $a+1$ exponent-two points, all
fixed. If $D_2$ is their divisor, the normalized cubic
algebra has eigensheaves
\[
 \pi_*\mathcal O_C
       =\mathcal O\oplus\mathcal L^{-1}
                          \oplus\mathcal L^{-2}(D_2).
                                                               \tag{H9.143}
\]
Indeed $u^3=z$ needs no adjustment, whereas $u^3=z^2$
has normalization $z=v^3,u=v^2$ and the third integral
basis element is $u^2/z=v$. These local calculations
give (H9.143), including infinity. Its quotient dual is
$\mathcal L\oplus\mathcal L^2(-D_2)=E_g$.
Both $\mathcal L$ and $D_2$ are constant over the moving
branch parameter, so this is an isomorphism of relative
bundles, not just an assertion about individual fibres.

In characteristic three use (H9.130), with $a+1$
simple poles (including zero) and $b$ double poles.
Their positions are fixed. Let $D_1$ have coefficient
one at every pole, and let $D_2'$ have coefficient one
at each simple pole and two at each double pole. Then
as modules the normalized algebra is
\[
 \pi_*\mathcal O_C
       =\mathcal O\oplus\mathcal O(-D_1)w
                          \oplus\mathcal O(-D_2')w^2.
                                                               \tag{H9.144}
\]
To check this, at a pole of order $d\in\{1,2\}$ the
valuations of the terms $a_j(t)w^j$, $0\le j\le2$,
are $3v(a_j)-jd$, which are pairwise distinct modulo
three. Thus their sum is integral exactly when each
term is integral, imposing $v(a_j)\ge\lceil jd/3\rceil$.
Away from the poles the Artin--Schreier algebra is
\'etale and has basis $1,w,w^2$. This proves (H9.144).
The same description holds over the parameter curve:
the explicit smooth normalizations in T23 show these
sections are regular, and the fibrewise basis test
makes the indicated map of locally free sheaves an
isomorphism. Since $\deg D_1=a+b+1$ and
$\deg D_2'=a+2b+1$, its quotient dual is again the
constant bundle $E_g$.

Choose these constant bundle isomorphisms. The
base-change-compatible equivalence (H9.136) gives
an actual morphism from the T23 parameter curve to
$U_g$ in either characteristic. In particular $U_g$
is a nonempty smooth irreducible open in $P_g$.

\paragraph{Why this parameter covers the generic point of the full space.}
We prove density, rather than assuming a Hurwitz-space
irreducibility theorem. Start with any cover whose
splitting type satisfies (H9.138)--(H9.140). If
$f_0-e\ge2$, the two sections $X$ and
$Y^{f_0-e-1}$ have no common zero and give an exact
sequence
\[
 0\longrightarrow\mathcal O(e)
   \xrightarrow{(X,\,Y^{f_0-e-1})}
   \mathcal O(e+1)\oplus\mathcal O(f_0-1)
   \longrightarrow\mathcal O(f_0)\longrightarrow0.
                                                               \tag{H9.145}
\]
The quotient is a line bundle of the displayed degree.
Let $\xi\in H^1(\mathcal O(e-f_0))$ be its extension
class. It is nonzero, since the middle bundle in
(H9.145) is not $\mathcal O(e)\oplus\mathcal O(f_0)$.
The family of classes $s\xi$ over $\mathbf A^1_s$
gives a bundle $\mathscr E$ on
$\mathbf P^1\times\mathbf A^1$: its fibre at zero is
$E$, while every nonzero fibre has splitting
$(e+1,f_0-1)$. A nonzero scalar multiple of an
extension class is obtained by an automorphism of
its sub-line bundle, which proves the latter assertion.

By (H9.140) and coherent cohomology and base change,
the direct image of
$\operatorname{Sym}^3\mathscr E\otimes\det\mathscr E^{-1}$
is locally free near zero and its fibre there is
$H^0(\mathbf P^1,\mathcal B_E)$. Thus the prescribed
section defining the original cover lifts to a
section over a parameter neighbourhood of zero.
For this exact consequence, see Grothendieck's
cohomology-and-base-change theorem, as stated with
its vanishing corollary in Conrad, \emph{Applications
of base change for coherent cohomology},
Theorem~1.1 and Corollary~1.2, p.~1.
Applying (H9.136) gives a finite flat family of covers.
Its special cover is smooth, connected, and
unramified above $\Sigma$, so shrinking around zero
preserves all three conditions. It follows that every
cover of splitting $(e,f_0)$ is a specialization of
covers of the more balanced type $(e+1,f_0-1)$.

Iteration proves that every geometric point of the
full smooth-cover stack is in the closure of the
balanced locus. That locus is open: at fixed total
degree $g+2$, balance is equivalent to
$H^0(E(-f_g-1))=0$, an open condition by
semicontinuity. It is the image of $U_g$, since a
choice of bundle isomorphism $E\simeq E_g$ presents
every balanced cover by (H9.136). In particular it
is irreducible. Its closure is the whole stack, so
the latter is irreducible and $U_g$ is a dominant
parameter for its generic point. The proof works
over every algebraically closed extension and in
every characteristic; it does not impose tame
ramification on general triple covers.

\paragraph{The actual universal trace kernel.}
Over $U_g$ form the smooth proper incidence family
\[
 \mathscr X=(\widetilde A\times U_g)
                    \times_{\mathbf P^1\times U_g}\mathscr C,
 \qquad p:\mathscr X\to U_g,
 \qquad Q:\mathscr X\to\widetilde A\times U_g.
                                                               \tag{H9.146}
\]
The map $Q$ is finite flat of degree three.
Smoothness follows over smooth pencil values by
smooth base change, and over $\Sigma$ because
the cover is \'etale there. Fix $\ell\ne\operatorname{char}k$.
On the lisse sheaf $R^np_*\mathbf Q_\ell$, the
correspondence
\[
                       P_{\rm new}=1-\frac13Q^*Q_*
                                                               \tag{H9.147}
\]
is an idempotent; $Q_*Q^*=3$. Its image is the
orthogonal complement of the constant pullback
$H^n(\widetilde A,\mathbf Q_\ell)$.

Restrict this sheaf and projector to the T23 curve
inside $U_g$ constructed above. Its image is exactly
the complete nontrivial deck contribution in T23,
including the fixed-ambient tensor contribution.
T23 proves this restricted sheaf has no nonzero
finite-orbit vector. The same is therefore true on
$U_g$: a vector with finite orbit under its
fundamental group would have finite orbit after
restriction to that curve, and the underlying
fibre vector remains nonzero. No surjectivity of
fundamental groups is required in this argument.

\paragraph{Algebraic classes, Chow averaging, and height.}
Put $F=\overline{k(U_g)}$ and take the geometric
generic cover. Every algebraic middle cohomology
class of $\mathcal X_F$ is defined over a finite
extension of $k(U_g)$ and hence has finite
monodromy orbit. Projecting it by (H9.147) gives
zero by the preceding paragraph. Thus for any
actual rational codimension-$m$ Chow cycle $Z$ on $\mathcal X_F$,
\[
           \left[Z-\frac13Q^*Q_*Z\right]=0
                       \quad\hbox{in }H^n(\mathcal X_F,\mathbf Q_\ell).
                                                               \tag{H9.148}
\]
This is a statement about its global cohomology;
there is no assertion that the cycle in brackets
vanishes in its Chow group.

All homologically trivial generic cycles are
admissible by H9-T12(iii), since the cover is
unramified over $\Sigma$. Choose admissible closures
$Z,W$ on $\mathcal X_F$. Pushforward gives
admissible closures $Q_*Z,Q_*W$ on $\widetilde A_F$;
this follows by testing against vertical cycles
and using finite-flat pullback and the projection
formula. Conversely their pullbacks are admissible
by the same test. Equation (H9.148) and the
intersection formula for admissible heights give
\[
 \begin{split}
 h_L(z,w)
  &=\frac1{9}\,(-1)^m
       \deg\bigl(Q^*Q_*Z\cdot Q^*Q_*W\bigr)\\
  &=\frac13\,(-1)^m
       \deg\bigl(Q_*Z\cdot Q_*W\bigr)
   =\frac13h_K(\operatorname{cor}z,\operatorname{cor}w).
 \end{split}                                                   \tag{H9.149}
\]
Here $(-1)^m$ is the common signed convention of
the preceding H9 height theorems. Since
$\operatorname{cor}\operatorname{res}=3$, corestriction
on rational homologically trivial Chow groups is
surjective. The asserted radical and height-quotient
identifications follow from (H9.149). Finally H9-T8
and accepted 032 give the stated fourfold sign.
\end{proof}

\paragraph{Sources and status.}
The all-base cubic equivalence and normalization
table are established background, directly read
in Wood, \emph{Algebra \& Number Theory} 5:8 (2011),
Theorem~2.1 and equation~(1), printed pp.~1073--1074
(publisher PDF pp.~6--7),
\url{https://msp.org/ant/2011/5-8/ant-v5-n8-p05-p.pdf}.
The elementary bundle-balancing, explicit splitting
of the T23 families, and resulting generic height
application are deductions supplied here. The
coherent base-change consequence was also directly
read in Conrad's cited handout,
\url{https://math.stanford.edu/~conrad/248BPage/handouts/cohom.pdf}.
The original base-change statements were also directly read in
EGA~III$_2$, Theorem~7.7.5(I--II), printed pp.~67--68
(continuous pp.~199--200), and Proposition~7.8.4, printed
pp.~73--74 (continuous pp.~205--206): on a reduced base,
constant fibrewise $h^0$ implies local freeness and tensor
base change; Theorem~7.7.5(II)(d) gives arbitrary scheme
base change. The required constancy here follows from
$H^1=0$ and Riemann--Roch.
\url{https://www.numdam.org/item/PMIHES_1963__17__5_0.pdf}.
All monodromy input is precisely the previously
checked T23 proof and its cited Katz theorems.
